%==============================================================
%============= Modèle de Travail de Maturité ================== %==============================================================

% On obtient le document final de la manière suivante :
% première compilation avec LaTeX
% seconde compilation avec LaTeX
% création de l'index : makeindex
% création de la biblographie : bibtex
% troisième compilation avec LaTeX
% éventuellement quatrième compilation avec LaTeX
% transformer le dvi e ps puis en PDF
% removed unused options: idxtotoc,bibtotoc,
\documentclass[10pt,a5paper,english,titlepage,twoside]{article}

%-------------------------
%------ PRÉAMBULE --------
%-------------------------

% Toutes les commandes techniques sont définies dans le fichier TM.sty
% Normalement, vous n'avez pas à y toucher
% Mais, il doit être présent
\usepackage{TM}

% Initialisation des paramètres du travail de maturité
% Il faut impérativement remplir ce fichier


% Définit le style des entête et pieds de page.
\pagestyle{fancy}

% Les ressources web sont définies dans une bibliographie à part.
%\newcites{web}{Webographie}



%---------------------------------
%------ DEBUT DU DOCUMENT --------
%---------------------------------

\begin{document}
%==============================================================
%====== FICHIER DE CONFIGURATION DU TRAVAIL DE MATURITÉ =======
%==============================================================
% Vous trouverez les instructions de compilation du fichier
% source principal main.tex dans la documentation se trouvant 
% dans le fichier DocumentationTMlatex.pdf

% Il est indispensable de remplir ce fichier !!!!!!!!!


%---------------------------------------------------------------
%-------------- La page de titre et page de garde --------------
%---------------------------------------------------------------

% Le titre du travail de maturité.
\newcommand{\worktitle}{The axioms of $n$-exact categories}
% Le sous-titre du travail. S'il n'y a pas de soustitre, laissez vide, pour un sous-titre sur plusieurs lignes, mettez un linebreak:
\newcommand{\worksubtitle}{}

% Texte au milieu de la page de titre.
\newcommand{\worktype}{Bachelor's thesis}

% L'auteur/ les auteurs
\newcommand{\theauthor}{Thomas Rosset}
\newcommand{\theadress}{Herman Krags Veg 45}
\newcommand{\lieu}{7050 Trondheim}
%si un seul auteur mettre partout des accolades vides {}
\newcommand{\theauthorB}{}
\newcommand{\theadressB}{}
\newcommand{\lieuB}{}

% La date de soumission du travail.
\newcommand{\workdatemonth}{Spring}
\newcommand{\workdateyear}{2025}

% Le thème du séminaire
\newcommand{\seminar}{}

% Définition du pied de page de titre en fonction de la langue (E nglish, F rancais or D eutsch)
% en décommentant la bonne ligne.
%\newcommand{\titlepagefooter}{\titlepagefooterE}
%\newcommand{\titlepagefooter}{\titlepagefooterD}
\newcommand{\titlepagefooter}{\titlepagefooterF}

\newcommand{\enseignant}{Johanne Haugland}
% Noter éventuellement la(les) branche(s) concernée(s)
\newcommand{\branche}{Abstract Algebra} % commenter cette ligne 
						% s'il y a une branche à préciser
						% et décommenter la suivante
%\newcommand{\branche}{nom de la branche}



%---------------------------------------------------------------
%------- Avant-propos, glossaire, remerciements, index ----------
%---------------------------------------------------------------

% Pour insérer un avant-propos : commentez décommentez Y es ou N o.
% Le fichier à remplir du texte désiré est pagesspeciales/avantpropos.tex
\newcommand{\unavantpropos}{Y}
%\newcommand{\unavantpropos}{N}

% Pour insérer un glossaire : commentez décommentez Y es ou N o.
% Le fichier à remplir du texte désiré est pagesspeciales/glossaire.tex
%\newcommand{\unglossaire}{Y}
\newcommand{\unglossaire}{N}

% Pour insérer des remerciements : commentez décommentez Y es ou N o.
% Le fichier à remplir du texte désiré est pagesspeciales/remerciements.tex
%\newcommand{\desremerciements}{Y}
\newcommand{\desremerciements}{Y}

% Pour insérer un index : commentez décommentez Y es ou N o.
% Il faut alors mettre les clés d'index avec \index{mot} et faire un make index.
%\newcommand{\unindex}{Y}
\newcommand{\unindex}{N}


%---------------------------------------------------------------
%------------ Chapitres et annexes -----------------------------
%---------------------------------------------------------------
% Le nombre de chapitres désirés
\newcommand{\nbchap}{3}
% Le nombre d'annexes désirées
\newcommand{\nbannexes}{}


%---------------------------------------------------------------
%--- Listes des figures, tables, listings, abréviations --------
%---------------------------------------------------------------
% Définition des noms des listes
%\renewcommand{\listfigurename}{Liste des figures} % Nom original = Table des figures
%\renewcommand{\listtablename}{Liste des tables} % Nom original = Liste des tableaux
%\renewcommand{\lstlistlistingname}{Liste des codes sources}  % Nom original = Listings

% Pour insérer une liste des figures : commentez décommentez Y es ou N o.
\newcommand{\unelistefig}{Y}
%\newcommand{\unelistefig}{N}

% Pour insérer une liste des tables : commentez décommentez Y es ou N o.
%\newcommand{\unelistetbl}{Y}
\newcommand{\unelistetbl}{N}

% Pour insérer une liste des listings : commentez décommentez Y es ou N o.
\newcommand{\unelistelst}{N}
%\newcommand{\unelistelst}{N}

% Pour insérer une liste des abréviations/notation : commentez décommentez Y es ou N o.
%\newcommand{\unelisteabreviations}{Y}
\newcommand{\unelisteabreviations}{N}


%---------------------------------------------------------------
%-------------- Objets flottants -------------------------------
%---------------------------------------------------------------
% Définitions des titres des figures, tables, listings...
% English: Figure, Table, Listing (default)
% German: Abbildung, Tabelle, Listing
% French: Figure, Table, Listing
\renewcommand{\figurename}{Figure}
\renewcommand{\tablename}{Table}


%---------------------------------------------------------------
%---------------- Figures --------------------------------------
%---------------------------------------------------------------
% Chemin vers le répertoire des figures. Vous pouvez ajouter d'autres chemins entre accolades
% à l'intérieur des accolades principales.
\graphicspath{{./images/}}

% Insertion de figures ; copier coller dans les chapitres en fonction des besoins
%\tmfigureB{NomFigureSansExtension}{Légende}{fig:votreLabel}	% Taille grande
%\tmfigureN{NomFigureSansExtension}{Légende}{fig:votreLabel}	% Taille normale
%\tmfigureS{NomFigureSansExtension}{Légende}{fig:votreLabel}	% Taille petite
%\tmfigureT{NomFigureSansExtension}{Légende}{fig:votreLabel}	% Taille très petite


%---------------------------------------------------------------
%--------------- Insertion de code : listings ------------------
%---------------------------------------------------------------
% Seul latex est appelé par défaut. On peut charger d'autres languages pour les listings
% selon le modèle de la ligne suivante. Voir le package listings pour les langages disponibles.
%\lstloadlanguages{HTML,PHP,TeX}
% Il faut ensuite placer [...,language=HTML,...] après l'appel \begin{lstlisting} pour définir
% le language à utiliser dans un script particulier


%---------------------------------------------------------------
%--------------- Version provisoire ----------------------------
%---------------------------------------------------------------

% Si vous travaillez avec une version provisoire, cette commande imprime un filigrane
% en haut à gauche de chaque page. Commentez la pour la version finale.
%\reviewtimetoday{\today}{* Draft Version *}


%---------------------------------------------------------------
%---------------- Le texte -------------------------------------
%---------------------------------------------------------------
% Réglage de l'indentation de la première ligne de chaque paragraphe
\parindent=0in		% pas d'indentation
%\parindent=0.1in	% petite indentation
%\parindent=0.2in	% moyenne indentation
%\parindent=0.3in	% grande indentation

% Si vous désirez commenter du code latex sur plusieurs lignes, utilisez la forme :
%\begin{comment} ... \end{comment}


%---------------------------------------------------------------
%----------------- Un index ------------------------------------
%---------------------------------------------------------------

% Il n'est pas nécessaire de réaliser un index. Par défaut, il n'y en a pas. Décommentez si vous en voulez un.
% Pour construire le fichier d'index idx, il doit être compilé avec makeindex en ligne de commande.
%\usepackage{makeidx} \makeindex

% Nettoyage des entête et pieds de page.
\pagestyle{empty}
\pagenumbering{gobble}
% Inclusion de la page de titre.
\begin{titlepage}
\begin{center}
	\seminar
	
	~\\[2 cm]
	
    
			

		
		 \begin{Huge}
        {\sf \worktitle \\}
    \end{Huge}
    \vspace{0.4cm}%
    \begin{LARGE} {\sf \worksubtitle} \end{LARGE}%

    % Si nécessaire, votre propre image
    %\vspace{1cm}
    %\includegraphics[height=1.5cm]{figures/projectLogo.eps}
   
    ~
    \\[2 cm]

    
     \begin{normalsize}
	{\textsc{\worktype}}\\ 
     \end{normalsize}
    ~
    \\[2  cm]
     \begin{normalsize}
	{\textsc{\theauthor}}\\
	\ifthenelse{\equal{\theauthorB}{}}{}{\textsc{\theauthorB} \\}
        \vspace{0.2cm}
     \end{normalsize}

    
    \vfill

    \titlepagefooter% modiification possible dans le fichier TM.sty

\end{center}

\end{titlepage}


% Inclusion de la page de garde.
\begin{center}
Bachelor's thesis completed\\
under supervision of\\
\enseignant

\bigskip
\ifthenelse{\equal{\branche}{}}{}{(\branche )}
\end{center}
\section*{Acknowledgements}
\addcontentsline{toc}{section}{Acknowledgements}

\section*{Abstract}
\setcounter{page}{3}

\pagenumbering{arabic}

% table des matières (fichier .toc)
\tableofcontents
%\renewcommand{\contentsname}{Table des matières} % Pour changer le titre


\addtolength{\parskip}{0.5\baselineskip} % espace inter-paragraphe (0.25 pour interligne simple)
%\remerciements
%\avantpropos
\pagestyle{fancy}
%------------------------------
% ----- CORPS DU DOCUMENT -----
%------------------------------

\section*{Introduction}
\addcontentsline{toc}{section}{Introduction}

\subsection{Prerequisites}
In order to understand this paper, the reader should have some knowledge about modern homological algebra. In particular, they should already be familiar with category theory, more specifically abelian category theory, as is presented for example in the lecture notes of MA3204 Homological Algebra \cite{homalg}.




\section*{Conventions and notations}\label{NOTATIONS}
\addcontentsline{toc}{section}{Conventions and notations}
Throughout this thesis, we let $n$ be a positive integer. We usually let $\cat{A}$ denote an additive category. When not stated otherwise, the word kernel (respectively cokernel) will refer to the kernel morphism (respectively cokernel morphism) instead of the object. For a morphism $f : A \xra{} B$, we call $A$ its domain and $B$ its codomain. For two objects $A, B \in \cat{A}$, we let $\cat{A}(A,B)$ denote the set of morphisms from $A$ to $B$. We will use $\ker{f}$ to refer to its kernel morphism and $\cok{f}$ for its cokernel morphism, whereas $\Ker{f}$ will refer to its kernel object and $\Cok{f}$ its cokernel object. For any object $A \in \cat{A}$, we denote by $1_A$ the identity morphism on $A$. We let $\Chlen{n}{\cat{A}}$ denote the set of chain complexes on $\cat{A}$ concentrated in degrees $\{0, \ldots, n\}$ and $\Ch{\cat{A}}$ denotes the set of all chain complexes on $\cat{A}$. We let $\text{H}(\cat{A})$ denote the homotopy category of $\cat{A}$



\leschapitres




%--------------------------
% ----- BIBLIOGRAPHIE -----
%--------------------------
%\renewcommand{\bibname}{Bibliographie}	 %Titre de la bibliographie en français
\bibliographystyle{plain}				% Style de la bibliographie
% Pour avoir toutes les références de vos fichiers bibtex dans la bibliographie,
% vous pouvez décommenter les lignes suivantes. Autrement, seules les références 
% cités apparaîtrons dans la bibliographie.
% \nocite{*}
%\nociteweb{*}
\bibliography{mainbib}
\addcontentsline{toc}{chapter}{Bibliography}

%-------------------
% ----- TABLES -----
%-------------------

% Liste des figures
%\listefig

% Liste des tables
\listetbl

% Ici vient la liste des listing/codes
% Définition du nom de la liste des listings
% Décommentez la ligne suivante et celle après les tirets si nécessaire

% Correction d'un petit bug de listing
\corrbug

% Liste des codes
\listelst 

% Liste des abreviations
\listeabreviations


%-------------------
%----- ANNEXES -----
%-------------------

\lesannexes

% Si nécessaire, l'index est placé ici.
\lindex

%\include{pagesspeciales/annexe1}
%\chapter*{Déclaration sur l'honneur}
\addcontentsline{toc}{chapter}{Déclaration sur l'honneur}


\begin{multicols}{2}
\theauthor\\
\theadress\\
\lieu\\
\columnbreak

\theauthorB\ \\
\theadressB\ \\
\lieuB\
\end{multicols}

\ifthenelse{\equal{\theauthorB}{}}
{Je certifie}
{Nous certifions}
que le travail

\begin{center}
{\Large\sf \worktitle}\\
{\large\sf \worksubtitle}
\end{center}


a été réalisé par
\ifthenelse{\equal{\theauthorB}{}}{moi}{nous}
conformément au \squote{Guide de travail} des collèges et aux \squote{Lignes directrices} de la DICS concernant la réalisation du \worktype.

\vspace{4cm}
Lieu, date: Bulle, 21.03.2020

%\includegraphics{images/annexes/signatures.png}





\end{document}
%--------------------------
%---- FIN DU DOCUMENT -----
%--------------------------
